\Askhsh{36355}{2}
\begin{ylh}{1}
    \item 4.6. Άθροισμα γωνιών τριγώνου
    \item 5.9. Μια ιδιότητα του ορθογώνιου τριγώνου.
\end{ylh}
% ===================================================================
%  Αρχείο: 36355.tex (Fragment)
%  Αυτόματη μετατροπή μέσω doc2latex.py
% ===================================================================

ΘΕΜΑ 2

Έστω ορθογώνιο τρίγωνο ΑΒΓ με $\widehat{Β} = 90^{ο}$ και Ζ το μέσο του ΑΓ. Με υποτείνουσα το ΑΓ κατασκευάζουμε ορθογώνιο ισοσκελές τρίγωνο ΑΔΓ με $\widehat{\Delta} = 90^{Ο}$ και ΔΑ = ΔΓ.

\begin{alist}
\item Να αποδείξετε ότι $ΒΖ = \Delta Ζ$. \monades{13}

\item Αν είναι $Α\widehat{\Gamma}Β = 30^{ο}$, τότε να υπολογίσετε τις γωνίες Β$\widehat{Α}$Δ και Β$\widehat{\Gamma}$Δ. \monades{12}


\begin{center}
\begin{tikzpicture}[scale=1]
\tkzDefPoint(0,0){B}
\tkzDefPoint(0,3){A}
\tkzDefPoint(4,0){Gamma}
\tkzDefMidPoint(A,Gamma) \tkzGetPoint{Z}
\tkzDefPointBy[rotation=center Z angle 90](A) \tkzGetPoint{Delta}

\draw[l3] (A)--(B)--(Gamma)--cycle;
\draw[l3] (A)--(Delta)--(Gamma);
\draw[l3] (B)--(Z);
\draw[l3] (Delta)--(Z);

\tkzDrawPoints(A,B,Gamma,Delta,Z)
\tkzLabelPoint[left](A){$A$}
\tkzLabelPoint[below left](B){$B$}
\tkzLabelPoint[right](Gamma){$\Gamma$}
\tkzLabelPoint[above right](Delta){$\Delta$}
\tkzLabelPoint[above](Z){$Z$}

\tkzMarkRightAngle(A,B,Gamma)
\tkzMarkRightAngle(A,Delta,Gamma)
\tkzMarkSegments[mark=s|](A,Z Z,Gamma B,Z Delta,Z)
\end{tikzpicture}
\end{center}
\end{alist}
